\fancyhead[L]{\small PASQUICH ÜBER SPHÄRISCHE ASTRONOMIE.}
\SourceLeaf{540}
\begin{center}\small Göttingische gelehrte Anzeigen. Stück 144. S. 1433..1436. 9. September 1811.\\[3pt]\rule{.72\linewidth}{.3pt}\end{center}

{\rightskip=0pt plus 1em \emph{Abriss der Grundlagen der sphärischen und rechnenden Astronomie von} \textsc{Joannes Pasquich}, Direktor der Sternwarte der Königlich Ungarischen Universität. \emph{Erster Teil. {\addfontfeatures{LetterSpace=10}Theoretische} Grundlagen der sphärischen und rechnenden Astronomie. Zweiter Teil, {\addfontfeatures{LetterSpace=10}praktische} Grundlagen der sphärischen und rechnenden Astronomie.} 1811. 160 und 166 Seiten in Quart. Wien, Schaumann u. Comp.\par}

Der Verfasser erklärt in der Vorrede, dass die Schrift aus astronomischen Vorlesungen an der Universität in Ofen entstanden sei. Sie richte sich allein an Anfänger ohne Vorkenntnisse und nicht an geübte Astronomen. Sie solle ihnen weniger eine vollständige Anleitung zu allen vielfältigen astronomischen Verfahren geben, als vielmehr ihr Interesse am Studium dieser dort bisher allzu sehr vernachlässigten Wissenschaft wecken und sie dazu ermuntern, damit sie später weiterführenden Unterricht erhalten können. Damit sei die sonst auffallende Dürftigkeit des Buches zu entschuldigen. Es ehrt den Verfasser, dass er dieses Ziel nicht durch ein oberflächliches Abschöpfen unterhaltsamer Ergebnisse zu erreichen sucht (was das ernsthafte Studium nicht fördert), sondern von einer gründlichen Vorbereitung in den Grundlagen ausgeht. Der erste, theoretische Teil behandelt in fünf Abschnitten die folgenden Themen. Der \emph{erste} Abschnitt trägt den Titel „Grundbegriff und Grundgesetze der täglichen Bewegung der Gestirne und der jährlichen Bewegung der Sonne“. Er gliedert sich in vier Kapitel. Sie erläutern die Grundbegriffe der Sphäre und der auf ihr gezogenen Kreise sowie die trigonometrischen Formeln, mit denen sich die Lage der Gestirne gegenüber diesen Kreisen bestimmen lässt. Der \emph{zweite} Abschnitt gliedert sich ebenfalls in vier Kapitel. Er behandelt die verschiedenen Arten von Zeit und den damit verbundenen Gebrauch der Rektaszensionen sowie die Unterschiede der Meridiane. Angesichts der Gründlichkeit, die sonst überall in der Darstellung herrscht, könnte man sich doch wünschen, dass hier der eine oder andere Punkt etwas vertieft worden wäre. So ist etwa der Begriff der mittleren Sonnenzeit als \emph{Zeitdauer} ganz klar, doch der Begriff der mittleren Sonnenzeit als \emph{Zeitpunkt} bleibt unscharf. Man sieht nicht, von wo aus man die mittlere und die wahre Bewegung der Sonne gleichzeitig beginnen lassen soll. Der \emph{dritte} Abschnitt handelt in drei Kapiteln von den Beziehungen zwischen der Bewegung der Sonne in der Ekliptik und ihren Rektaszensionen, Deklinationen und Positionswinkeln sowie von den Höhen, Stundenwinkeln und Azimuten der Gestirne. Der \emph{vierte} Abschnitt behandelt in vier Kapiteln die Gestalt und Größe der Erde, die Parallaxe und die astronomische Lichtbrechung. Im \emph{fünften} Abschnitt vermittelt der Verfasser in fünf Kapiteln die Grundlagen der elliptischen Bewegung der Planeten und der davon abhängigen Erscheinungen sowie die Wirkungen der Präzession, Nutation und Aberration. Damit endet der theoretische Teil. Aus dieser Inhaltsübersicht sieht man, dass es dem Verfasser mehr darum ging, eine Sammlung brauchbarer Rechenformeln bereitzustellen, als die Darstellung so zu ordnen, dass man ohne mündliche Nachhilfe zu einer lehrreichen und befriedigenden Einsicht darin gelangen könnte, wie die Wahrheiten der Astronomie gefunden wurden oder gefunden werden könnten. Dieser Nachhilfe muss also vieles überlassen bleiben.

Der \emph{zweite Teil} des Werks ist keineswegs als vollständiges Handbuch der rechnenden Astronomie anzusehen. Er beschränkt sich auf die praktischen Anleitungen, die ein beobachtender Astronom wohl am häufigsten braucht. Er enthält drei Abschnitte. Der \emph{erste} erläutert für die Praxis, wie astronomische Tafeln im Allgemeinen und Sonnentafeln im Besonderen aufgebaut sind, wie man astronomische\SourceLeaf{541} Ephemeriden und Fixsternverzeichnisse gebraucht und wie man Nutation und Aberration berechnet. Der \emph{zweite} Abschnitt lehrt den Vergleich der verschiedenen astronomischen Zeiten miteinander und die Zeitbestimmung aus beobachteten Kulminationen und korrespondierenden Höhen. Der \emph{dritte} Abschnitt zeigt vor allem, wie sich beobachtete Zirkummeridianhöhen zur Breitenbestimmung verwenden lassen, und beschreibt, wie die Stellungen der Himmelskörper aus den Beobachtungen abzuleiten sind.

Zum oben genannten Werk gehört noch ein Anhang zu Joannes Pasquichs Abriss der praktischen Grundlagen der sphärischen und rechnenden Astronomie mit Hilfstafeln. Wien bei Schaumburg et Soc. 1811. 42 Seiten in Quart. Es sind dies die abgekürzten Sonnentafeln von \textsc{Delambre}, die allgemeine Tafel für die Mittagskorrektur in der üblichen Anordnung, eine Hilfstafel zur Reduktion der außerhalb des Meridians gemessenen Zenitdistanzen, eine weitere Hilfstafel zur Reduktion der Zenitdistanzen des Polarsterns, eine Tafel zur Bestimmung der Präzession in Rektaszension und Deklination, Tafeln für Refraktion, Aberration und Nutation und der neueste Katalog von \textsc{Piazzi} mit 121 Fundamentalsternen. Der Rezensent hätte sich gewünscht, dass einigen dieser Tafeln, deren Anordnung für sich genommen nicht verständlich ist (etwa den Tafeln zur Reduktion der Zenitdistanzen), eine Erklärung beigefügt wäre, da manche die Tafeln vielleicht ohne den Abriss erwerben möchten.
